% GENERATED FILE. Edit canonical lessons, then run npm run generate:books.
% canonical-source-hash: fnv1a64:fa72a0db2aa2b99c
% canonical-lessons: UR-C46-ustad, UR-C46-khana, UR-C46-pina, UR-C46-dekhna, UR-C46-sunna

\chapter{Teacher, To eat, To drink, To see, To hear}
\label{ch:ur-teacher-to-eat-to-drink-to-see-to-hear}
\hlchaptermodality{46}

\begin{chapteropening}
\textbf{By the end of this chapter:} \emph{I can say a teacher, to eat, to drink, to see, to look and to hear, to listen.}

Complete the last lesson of chapter 46: I can say a teacher, to eat, to drink, to see, to look and to hear, to listen.
\end{chapteropening}

\section[\texorpdfstring{\ur{استاد} (ustād)}{استاد (ustād)}]{\ur{استاد} (ustād) — a teacher}

\label{lesson:UR-C46-ustad}



\begin{quote}
Before the new one: say the Urdu for a neighbour, then the Urdu for a guest.
\end{quote}

\subsection*{You'll want to know: \ur{استاد}}
\textbf{\ur{استاد}} — \emph{ustād} — \textquotedblleft{}a teacher\textquotedblright{}.

\begin{cousinweb}
\textbf{\ur{استاد}} (\emph{ustād}) is a loan, usually traced to Persian \emph{ustād}, \textquotedblleft{}master, teacher\textquotedblright{}.
\end{cousinweb}

\begin{grammarlens}[title={What you've built}]
One more word for this chapter.
\end{grammarlens}

\subsection*{Your turn}
\begin{itemize}
\raggedright
  \item \emph{Say it:} \emph{ustād}
  \item \emph{Say it:} \emph{ustād}, once more
  \item \emph{Say it:} say \emph{ustād}
  \item \emph{Recall:} say the Urdu for a neighbour, then the Urdu for a guest, then say \emph{ustād} again
\end{itemize}

\begin{tcolorbox}[breakable,colback=teal!4,colframe=teal!35!black,title={Before you move on}]
What does \ur{استاد} mean? (\textquotedblleft{}A teacher\textquotedblright{}.) Say it once more.
\end{tcolorbox}
\section[\texorpdfstring{\ur{کھانا} (khānā)}{کھانا (khānā)}]{\ur{کھانا} (khānā) — to eat; food}

\label{lesson:UR-C46-khana}



\begin{quote}
Before the new one: say the Urdu for a guest, then the Urdu for a teacher.
\end{quote}

\subsection*{You'll want to know: \ur{کھانا}}
\textbf{\ur{کھانا}} — \emph{khānā} — \textquotedblleft{}to eat; food\textquotedblright{}.

\begin{cousinweb}
\textbf{\ur{کھانا}} (\emph{khānā}) is two inherited words that met in one shape. The verb \textquotedblleft{}to eat\textquotedblright{} comes from Sanskrit \emph{khādati}, \textquotedblleft{}chews, eats\textquotedblright{} (CDIAL 3865), and the noun \textquotedblleft{}food\textquotedblright{} from \emph{khādana}, \textquotedblleft{}eating, food\textquotedblright{} (CDIAL 3867).
\end{cousinweb}

\begin{grammarlens}[title={What you've built}]
One more word for this chapter.
\end{grammarlens}

\subsection*{Your turn}
\begin{itemize}
\raggedright
  \item \emph{Say it:} \emph{khānā}
  \item \emph{Say it:} \emph{khānā}, once more
  \item \emph{Say it:} say \emph{khānā}
  \item \emph{Recall:} say the Urdu for a guest, then the Urdu for a teacher, then say \emph{khānā} again
\end{itemize}

\begin{tcolorbox}[breakable,colback=teal!4,colframe=teal!35!black,title={Before you move on}]
What does \ur{کھانا} mean? (\textquotedblleft{}To eat; food\textquotedblright{}.) Say it once more.
\end{tcolorbox}
\section[\texorpdfstring{\ur{پینا} (pīnā)}{پینا (pīnā)}]{\ur{پینا} (pīnā) — to drink}

\label{lesson:UR-C46-pina}



\begin{quote}
Before the new one: say the Urdu for a teacher, then the Urdu for to eat; food.
\end{quote}

\subsection*{You'll want to know: \ur{پینا}}
\textbf{\ur{پینا}} — \emph{pīnā} — \textquotedblleft{}to drink\textquotedblright{}.

\begin{cousinweb}
\textbf{\ur{پینا}} (\emph{pīnā}) is a tadbhava, a word inherited from Sanskrit and reshaped by sound change. It comes from \emph{pibati}, \textquotedblleft{}drinks\textquotedblright{} (CDIAL 8209).
\end{cousinweb}

\begin{grammarlens}[title={What you've built}]
One more word for this chapter.
\end{grammarlens}

\subsection*{Your turn}
\begin{itemize}
\raggedright
  \item \emph{Say it:} \emph{pīnā}
  \item \emph{Say it:} \emph{pīnā}, once more
  \item \emph{Say it:} say \emph{pīnā}
  \item \emph{Recall:} say the Urdu for a teacher, then the Urdu for to eat; food, then say \emph{pīnā} again
\end{itemize}

\begin{tcolorbox}[breakable,colback=teal!4,colframe=teal!35!black,title={Before you move on}]
What does \ur{پینا} mean? (\textquotedblleft{}To drink\textquotedblright{}.) Say it once more.
\end{tcolorbox}
\section[\texorpdfstring{\ur{دیکھنا} (dekhnā)}{دیکھنا (dekhnā)}]{\ur{دیکھنا} (dekhnā) — to see, to look}

\label{lesson:UR-C46-dekhna}



\begin{quote}
Before the new one: say the Urdu for to eat; food, then the Urdu for to drink.
\end{quote}

\subsection*{You'll want to know: \ur{دیکھنا}}
\textbf{\ur{دیکھنا}} — \emph{dekhnā} — \textquotedblleft{}to see, to look\textquotedblright{}.

\begin{cousinweb}
\textbf{\ur{دیکھنا}} (\emph{dekhnā}) is a tadbhava, a word inherited and reshaped by sound change. It comes from \emph{*dekṣati}, \textquotedblleft{}sees\textquotedblright{}, a form CDIAL reconstructs from its descendants (CDIAL 6507). It is a variant CDIAL reconstructs beside \emph{*dṛkṣati}.
\end{cousinweb}

\begin{grammarlens}[title={What you've built}]
One more word for this chapter.
\end{grammarlens}

\subsection*{Your turn}
\begin{itemize}
\raggedright
  \item \emph{Say it:} \emph{dekhnā}
  \item \emph{Say it:} \emph{dekhnā}, once more
  \item \emph{Say it:} say \emph{dekhnā}
  \item \emph{Recall:} say the Urdu for to eat; food, then the Urdu for to drink, then say \emph{dekhnā} again
\end{itemize}

\begin{tcolorbox}[breakable,colback=teal!4,colframe=teal!35!black,title={Before you move on}]
What does \ur{دیکھنا} mean? (\textquotedblleft{}To see, to look\textquotedblright{}.) Say it once more.
\end{tcolorbox}
\section[\texorpdfstring{\ur{سننا} (sunnā)}{سننا (sunnā)}]{\ur{سننا} (sunnā) — to hear, to listen}

\label{lesson:UR-C46-sunna}



\begin{quote}
Before the new one: say the Urdu for to drink, then the Urdu for to see.
\end{quote}

\subsection*{You'll want to know: \ur{سننا}}
\textbf{\ur{سننا}} — \emph{sunnā} — \textquotedblleft{}to hear, to listen\textquotedblright{}.

\begin{cousinweb}
\textbf{\ur{سننا}} (\emph{sunnā}) is a tadbhava, a word inherited from Sanskrit and reshaped by sound change. It comes from \emph{śṛṇoti}, \textquotedblleft{}hears\textquotedblright{} (CDIAL 12598).
\end{cousinweb}

\begin{grammarlens}[title={What you've built}]
One more word for this chapter.
\end{grammarlens}

\subsection*{Your turn}
\begin{itemize}
\raggedright
  \item \emph{Say it:} \emph{sunnā}
  \item \emph{Say it:} \emph{sunnā}, once more
  \item \emph{Say it:} say \emph{sunnā}
  \item \emph{Recall:} say the Urdu for to drink, then the Urdu for to see, then say \emph{sunnā} again
\end{itemize}

\begin{tcolorbox}[breakable,colback=teal!4,colframe=teal!35!black,title={Before you move on}]
What does \ur{سننا} mean? (\textquotedblleft{}To hear, to listen\textquotedblright{}.) Say it once more.
\end{tcolorbox}
